\documentclass[11pt]{article}
\usepackage[margin=1in]{geometry}
\usepackage{graphicx}
\usepackage{booktabs}
\usepackage{siunitx}
\usepackage{amsmath}
\usepackage{hyperref}
\usepackage{listings}
\usepackage{caption}

\title{Keel Plate Line-Heating Report}
\author{Generated from PlateLineheating Results}
\date{\today}

\lstset{
  basicstyle=\ttfamily\small,
  breaklines=true,
  frame=single
}

\begin{document}
\maketitle

\section{Input IGES Geometry}
The IGES surface (\texttt{IGESFiles/keel plate 0-5.igs}) is a NURBS surface.
The extracted geometric measures used in the solver are summarized in Table~\ref{tab:iges}.

\begin{table}[h!]
\centering
\caption{IGES-derived curvature metrics}
\label{tab:iges}
\begin{tabular}{lrr}
\toprule
Metric & Value & Units \\
\midrule
Z-range (deflection range) & 98.00 & mm \\
Mean Z & 750.04 & mm \\
Target radius $R$ & 11479.59 & mm \\
Target curvature $\kappa=1/R$ & $8.711\times10^{-5}$ & 1/mm \\
\bottomrule
\end{tabular}
\end{table}

\section{FEM Solver Mathematics}
The thermo-mechanical model follows a coupled line-heating formulation:

\subsection*{Heat conduction}
\begin{equation}
\rho c_p \frac{\partial T}{\partial t} - \nabla \cdot (k \nabla T) = Q(x,y,z,t)
\end{equation}
where $Q$ is the moving Gaussian heat source along each heating line. Convection is applied on surfaces:
\begin{equation}
-k\nabla T \cdot \mathbf{n} = h_{\mathrm{conv}}(T - T_{\infty}).
\end{equation}

\subsection*{Mechanical equilibrium}
\begin{equation}
\nabla \cdot \boldsymbol{\sigma} + \mathbf{b} = 0,
\end{equation}
with linear elasticity and thermal/inherent strains
\begin{equation}
\boldsymbol{\varepsilon} = \boldsymbol{\varepsilon}^{\mathrm{el}} + \boldsymbol{\varepsilon}^{\mathrm{th}} + \boldsymbol{\varepsilon}^{\mathrm{inh}}.
\end{equation}
Thermal strain $\boldsymbol{\varepsilon}^{\mathrm{th}} = \alpha (T - T_{\mathrm{ref}})\mathbf{I}$,
and the inherent strain $\boldsymbol{\varepsilon}^{\mathrm{inh}}$ is applied in a band around the heating lines.

\section{Line Pattern and Heat Input}
The final optimized line pattern uses five parallel lines across the plate width:

\begin{itemize}
  \item Number of lines: 5
  \item Line positions (y): 166.67, 333.33, 500.00, 666.67, 833.33 mm
  \item Line energy: 580 J/mm
  \item Heat flux peak: $q_0 = 5800.01\ \mathrm{W/mm^2}$
  \item Velocity: 10 mm/s
  \item Passes: 1
\end{itemize}

Figure~\ref{fig:lines} shows the heating line profiles on the plate.

\begin{figure}[h!]
  \centering
  \includegraphics[width=0.9\textwidth]{../results/keel_plate_opt_loop/iter_03/heating_line_profiles.png}
  \caption{Heating line pattern and heat flux profiles (optimized run).}
  \label{fig:lines}
\end{figure}

\section{FEM Results (Optimized Run)}
The optimized simulation (iteration 3) achieves the target curvature within 0.3\%.

\begin{table}[h!]
\centering
\caption{IGES target vs FEM results (optimized run)}
\label{tab:comparison}
\begin{tabular}{lrrr}
\toprule
Metric & IGES Target & FEM Result & Units \\
\midrule
Radius $R$ (transverse) & 11479.59 & 11510.89 & mm \\
Curvature $\kappa$ (transverse) & $8.711\times10^{-5}$ & $8.687\times10^{-5}$ & 1/mm \\
Deflection range & 98.000 & 19.306 & mm \\
Max deflection & -- & 5.318 & mm \\
Min deflection & -- & -13.988 & mm \\
\bottomrule
\end{tabular}
\end{table}

Figure~\ref{fig:deflection} shows the deformed shape from the FEM solution.

\begin{figure}[h!]
  \centering
  \includegraphics[width=0.9\textwidth]{../results/keel_plate_opt_loop/iter_03/deflection_3d.png}
  \caption{FEM deflection field (optimized run).}
  \label{fig:deflection}
\end{figure}

\begin{figure}[h!]
  \centering
  \includegraphics[width=0.9\textwidth]{../results/keel_plate_opt_loop/iter_03/deflection_3d_overlay.png}
  \caption{3D deformation overlay (VTK-derived).}
  \label{fig:deflection_overlay}
\end{figure}

\begin{figure}[h!]
  \centering
  \includegraphics[width=0.9\textwidth]{../results/keel_plate_opt_loop/iter_03/deflection_3d_unscaled_vtk.png}
  \caption{Unscaled deformation from VTK (deflection surface).}
  \label{fig:deflection_unscaled}
\end{figure}

\section{Input Files}
The final FEM inputs and the optimization script are included below.

\subsection*{FEM Input JSON}
\lstinputlisting{../config/keel_plate_opt_loop_final.json}

\subsection*{Python Optimization Script}
\lstinputlisting{../run_keel_optimize_loop_report.py}

\section{Deflection Data}
The node-wise deflection array is stored in:
\begin{itemize}
  \item \texttt{results/keel\_plate\_opt\_loop/iter\_03/displacement.npy}
\end{itemize}
and the summary statistics are in:
\begin{itemize}
  \item \texttt{results/keel\_plate\_opt\_loop/iter\_03/summary.json}
  \item \texttt{results/keel\_plate\_opt\_loop/iter\_03/curvature\_comparison.json}
\end{itemize}

\end{document}